\section{Monte Carlo 用随机样本逼近期望}
\label{sec:13-Machine-Learning/08-Sampling-and-Inference/01}

电网调度组问你一个听起来很简单的问题：明年冬天出现大面积停电的损失期望是多少。你已经有了一个还算像样的模型——天气、负载、几十台设备的老化状态构成随机向量 \(X\)，联合密度 \(p(x)\) 写得清清楚楚，损失 \(h(X)\) 由一段仿真代码算出。目标量只有一行：

\[
\mu=\E_p[h(X)]
=\int h(x)p(x)\,dx.
\]

写得出，算不出。这是个几十维的积分，\(h\) 由分支逻辑拼成，既非线性也不光滑，解析求解无从下手。铺规则网格呢？每维取 \(m\) 个点就是 \(m^d\) 次仿真，\(d=20\)、\(m=10\) 已经是 \(10^{20}\) 次函数求值。整个单元要处理的困境从这里开始：目标量的表达式是完整的，通往数值的那条路却被维数堵死了。

Monte Carlo 换了一条路。不铺网格，随机撒点：从 \(p\) 独立抽 \(X_1,\ldots,X_n\)，直接取平均，

\[
\widehat\mu_n
=\frac1n\sum_{i=1}^nh(X_i).
\]

积分变成了平均，而平均的误差按 \(1/\sqrt n\) 缩小，与维数 \(d\) 无关。维数只通过 \(h(X)\) 的方差间接起作用，不改变误差随样本量衰减的速率。你可以在千维空间里撒点，只要 \(h\) 的波动不过分，几百个样本就够给出一个可用的数。

合法性来自两条经典结果。只要 \(\E\abs{h(X)}<\infty\)，大数定律保证 \(\widehat\mu_n\to\mu\)；记 \(\sigma_h^2=\Var(h(X))\)，若 \(\sigma_h^2<\infty\)，则 \(\Var(\widehat\mu_n)=\sigma_h^2/n\)，中心极限定理进一步给出

\[
\sqrt n\,(\widehat\mu_n-\mu)
\overset{d}{\longrightarrow}
\mathcal N(0,\sigma_h^2).
\]

把 \(\sigma_h\) 换成样本标准差 \(s_h\)，就得到 Monte Carlo 标准误 \(s_h/\sqrt n\)。这是整套方法的工作马：能从 \(p\) 采样、能算 \(h\)，就能交出一个带误差棒的估计。

\begin{center}
\begin{tikzpicture}[>=stealth]
\draw[gray!60] (0,0) -- (9.4,0);
\draw[->] (0,-3.1) -- (0,3.1) node[above]{\(\widehat\mu_n-\mu\)};
\draw[->] (0,0) -- (9.6,0) node[right]{\(n\)};
\draw[dashed] plot[domain=0.7:9.2,samples=90] (\x,{2.4/sqrt(\x)});
\draw[dashed] plot[domain=0.7:9.2,samples=90] (\x,{-2.4/sqrt(\x)});
\node at (7.4,1.45) {\(+1.96\,\sigma_h/\sqrt n\)};
\node at (7.4,-1.45) {\(-1.96\,\sigma_h/\sqrt n\)};
\draw[thick] plot[smooth] coordinates
  {(0.7,2.5) (1.1,-1.5) (1.6,1.4) (2.1,0.2) (2.7,-1.2)
   (3.3,0.8) (4.0,-0.6) (4.8,0.5) (5.7,-0.55) (6.6,0.35)
   (7.5,-0.3) (8.4,0.25) (9.2,-0.18)};
\draw[gray!60,dashed] (4,-3.1) -- (4,3.1);
\node[gray] at (4.75,-2.75) {四倍样本};
\draw[gray!60,dashed] (8,-3.1) -- (8,3.1);
\end{tikzpicture}
\end{center}

\emph{误差带按 \(n^{-1/2}\) 收窄，估计值在带内游走；把误差减半要把样本量翻两番。}

\subsection{无偏说的是长期重复，不是这一次运行}

\(\widehat\mu_n\) 无偏，意思是把整套模拟独立重复无穷多次再取平均才对。它不承诺某一次运行已经足够接近。图里那条游走的曲线正是一次运行：它可以在很长一段区间里稳定地偏在零线一侧，看上去像收敛到了别的地方。判断收没收敛，要看它有没有留在收窄的带子里，而不是看它最近几百步平不平。误差按 \(n^{-1/2}\) 缩小，想让标准误减半得付出四倍算力——图中两条竖线之间就是这笔账。多打印几位小数不增加任何信息，只是把噪声印得更漂亮。

除了损失期望，调度组多半还关心越限概率：

\[
P(h(X)>c)
=\E\bigl[\ind\set{h(X)>c}\bigr].
\]

这仍然是一个期望，只不过被积函数换成了指示函数。麻烦在于事件罕见的时候。若真实概率是 \(10^{-5}\)，跑一百万次仿真可能一次越限都撞不上，估计值直接给出零。零不表示风险消失了，只表示采样分布把几乎全部预算花在了普通情形上，极少访问那个要命的尾部。此时标准误公式也失效：样本全零则 \(s_h=0\)，标准误估成零，而真实误差可能很大。一个报告出 \(0\pm0\) 的罕见事件估计，是最危险的那种输出，它的形式完美而内容为空。

\subsection{方差缩减比盲目加样本便宜得多}

如果 \(n=10^4\) 时方差仍嫌大，把 \(n\) 提到 \(10^6\) 要一百倍算力。更划算的问法是：同样的样本量下，能不能让方差变小。

重要性采样改从另一个分布 \(q\) 抽样：

\[
\E_p[h(X)]
=\E_q\!\left[h(X)\frac{p(X)}{q(X)}\right].
\]

把右端按定义展开，\(q(x)\) 在分子分母上约掉，恢复成 \(\int h(x)p(x)\,dx\)，恒等式就这么一步。但换了采样分布，方差可以差出几个数量级：让 \(q\) 把概率质量压到对积分贡献大的区域，每个样本携带的信息就更多。罕见事件正是它的主场——把 \(q\) 偏向越限区，一万个样本里可能有几百个落在尾部，而不是一个都没有。

代价藏在权重 \(w(x)=p(x)/q(x)\) 里。第一条硬约束是支撑覆盖：\(p(x)>0\) 而 \(q(x)=0\) 的地方永远采不到，估计带上无法修复的偏差。即便支撑覆盖了，第二条更隐蔽——\(q\) 的尾部比 \(p\) 薄时，权重在尾部按指数发散，权重的方差可以是无穷的，此时中心极限定理的前提根本不成立。可观察的诊断量是有效样本量

\[
n_{\mathrm{eff}}
=\frac{\bigl(\sum_i w_i\bigr)^2}{\sum_i w_i^2},
\]

它在权重全相等时等于 \(n\)，在一个权重独大时趋于 \(1\)。表面样本量十万而 \(n_{\mathrm{eff}}\) 只有十几，说明这个平均实际上由十几个样本主导，大数定律远未生效。跑完之后至少看三件事：\(n_{\mathrm{eff}}/n\) 的比值、最大权重占总权重的份额、以及 \(\log w\) 的直方图右尾是否拖出长长一条。三者中任何一个报警，都不该把区间当真。

控制变量走的是另一条路。找一个与 \(h(X)\) 相关、期望 \(\E g\) 已知的量 \(g(X)\)，构造

\[
\widehat\mu_{\mathrm{cv}}
=\frac1n\sum_i\bigl[h(X_i)-c\,(g(X_i)-\E g)\bigr].
\]

想法很朴素：\(h\) 和 \(g\) 同涨同跌，那么 \(h\) 的波动里有一部分可以用 \(g\) 的波动来解释，减掉它。修正项期望为零，不引入偏差。方差关于 \(c\) 是二次的，

\[
n\Var(\widehat\mu_{\mathrm{cv}})
=\sigma_h^2-2c\Cov(h,g)+c^2\sigma_g^2,
\]

在 \(c^{*}=\Cov(h,g)/\sigma_g^2\) 处取到最小值 \(\sigma_h^2(1-\rho^2)\)，\(\rho\) 是 \(h\) 与 \(g\) 的相关系数。相关越强，方差降得越多；相关系数 \(0.9\) 就能把方差砍掉八成。分层抽样保证每个子群按配额被覆盖，对偶变量用负相关的样本对互相抵消波动。这些手法的共同逻辑是同一条：保持目标期望不变，把采样预算挪到分布中贡献大、波动大的地方去。想清楚积分的质量集中在哪里，比盲目加样本有效得多。

\subsection{方差无穷时误差棒不再有意义}

前面所有标准误公式都挂在 \(\sigma_h^2<\infty\) 这个条件上。若 \(h(X)\) 重尾到方差发散——某些金融损失、极端天气赔付确实如此——那么 \(\widehat\mu_n\pm1.96\,s_h/\sqrt n\) 不再有通常的覆盖概率，样本均值可能收敛到某个 \(\alpha<2\) 的稳定分布而不是正态，收敛速度也慢得多。重要性采样权重的重尾同样致命，且更容易被忽略，因为它藏在中间量里而不是最终输出里。

遇到可疑情形，先做批次诊断：把总预算切成若干独立批，看批间标准差是否随批容量按 \(n^{-1/2}\) 缩小。若明显更慢，或者批间散布被个别批主导，那就是二阶矩不存在的信号。补救手段包括重新设计提议分布、对权重做截断或 Pareto 平滑、改用分位数一类不依赖二阶矩的目标量。

一份可信的模拟报告不止一行估计值加一行区间。它记录随机种子、样本数、点估计与 Monte Carlo 标准误，附上多批独立运行的误差-样本量曲线，并在存在解析解的简化版本上先核对过实现。最后还有一条边界必须写在报告里：Monte Carlo 误差与模型误差是两件事。样本量趋于无穷，只会更精确地算出当前模型下的答案——当前对天气分布的假设、对损失函数的近似、对参数的估计。假设错了，撒再多点也修不回来。而当你连从 \(p\) 独立采样都做不到时——后验分布通常就是这样，密度只知道差一个常数——下一节要构造的那条链，是绕开这道障碍的另一条路。
